\section{Reflection as a possible explanation of cumulative positivity}
\label{sec:storage}

\subsection{The identity the positive pairing must supply}

The second construction target is a positive Hilbert space
$\mathcal H_{\omega,L}$ and a linear amplitude map
$\mathcal B_{\omega,L}:L^2(I_L)\to\mathcal H_{\omega,L}$ such that
\begin{equation}\label{eq:norm-hypothesis}
 \boxed{\norm f_{L^2(I_L)}^2
 =\norm{\mathcal V_{\omega,L}f}_{L^2(I_L)}^2
 +\norm{\mathcal B_{\omega,L}f}_{\mathcal H_{\omega,L}}^2.}
\end{equation}
This is a proposed physical balance law. Once the transfer
matching has been proved, it implies
\begin{equation}
 D_{\omega,L}=\mathcal B_{\omega,L}^*\mathcal B_{\omega,L}
 \succeq0.
\end{equation}
A proof first on a dense core, with bounded extensions, is enough.
The initial normalization forces $\mathcal B_{0,L}=0$ if the
identity extends to zero shift.

The input and output norms in \eqref{eq:norm-hypothesis} are the
fixed, unweighted arithmetic norms. A parameter-dependent physical
metric requires an explicit isometric identification and the extra
terms caused by its variation. Defining $\mathcal B=D^{1/2}$ after
assuming $D\succeq0$ would restate the desired conclusion; the
amplitudes must be specified independently. A whole-line unitary
model available only under RH also cannot be the starting proof.

\subsection{What ordinary reflection currently establishes}

For tangent map $M$ and time reflection
$R=\diag(-1,1,1,1)$, the adjoint branch uses $N=-MR$ and reversed
path ordering. If $F_i$ is the Wilson-dressed endpoint vector at a
common reference, the proposed network is
\[
 K_{ij}=\vev{\Theta_0F_i\,F_j}.
\]
A regulator with the appropriate reference color sector and
Osterwalder--Schrader inequality would make this a Gram matrix.
The interacting regulator, operator domain and positivity-preserving
renormalization are assumptions, not established consequences of
gauge invariance. The full construction and its scope are in
\cite[Section~S7]{Supplement}.

At free endpoint order the time-ray kernel is unconditionally
positive:
\begin{equation}
 K^{(0)}(r,s)=\frac1{r+s}
 =\int_0^\infty e^{-\lambda r}e^{-\lambda s}\dd\lambda,
 \qquad r,s>0.
\end{equation}
After radial weights it gives
$G_o(u)=1/(2\cosh(u/2))$.
This is only the opposite-ray part of
$\nGamma(|u|)=(G_s(u)+G_o(u))/2$, where
$G_s(u)=1/(2\sinh(|u|/2))$.
Moreover, a positive kernel of states has not been identified
with the defect $D_{\omega,L}$.
It needs a specified shift-dependent amplitude and the exact balance
\eqref{eq:norm-hypothesis}. The unmodified free kernel does not
satisfy the required vanishing at zero shift.

Internal rotations that align the two scalar maps need not
preserve this positive metric; explicit endpoint-doublet and odd-H
counterexamples are retained in the companion. The positive
ordinary adjoint is therefore a substantive part of the proposal.

\subsection{Why the target is the accumulated norm}

Equation~\eqref{eq:cumulative} integrates the energy on the moving
input $V_{s,L}f$. A reflected realization might explain the resulting
positive norm without factoring each instantaneous $Q_{s,L}$.
The shifted-zeta storage calculation records a useful distinction:
at $L=\log7$ and $\omega=10^{-11}$, one exact polynomial input has
a negative instantaneous quotient of about $-2.998\times10^{-27}$
but a positive cumulative quotient of about
$8.7170\times10^{-38}$
\cite[Appendix~A]{StorageDepth}.
These inherited enclosures concern one input, not an all-input
contraction bound. Nor do they show negativity of the actual
evolving integrand, since $Q_{\omega,L}[f]$ and
$Q_{\omega,L}[V_{\omega,L}f]$ use different inputs.

The example motivates allowing accumulated storage to carry the
positivity argument. It does not establish the physical balance
law or justify reversing a forward dissipation inequality.
